\documentclass[11pt,a4paper]{article}

\usepackage[T1]{fontenc}
\usepackage[utf8]{inputenc}
\usepackage[polish]{babel}
\usepackage[margin=2.5cm]{geometry}

\setlength{\parindent}{0pt}
\setlength{\parskip}{6pt}

\begin{document}

\begin{center}
    {\Large\textbf{Drzewo BST}}\\[4pt]
    Dokumentacja metod usuwania
\end{center}

\section{Reprezentacja drzewa}

Drzewo jest zaimplementowane w klasie \texttt{BST}.
Węzły są przechowywane w kontenerze \texttt{std::vector<Node>}.
Każdy węzeł zawiera wartość całkowitą oraz indeksy lewego
i prawego dziecka. Indeks równy \texttt{-1} oznacza brak dziecka.
Pole \texttt{rootIndex} wskazuje korzeń drzewa.

\section{Usuwanie pojedynczego węzła}

Metoda \texttt{removeNode(int value)} wyszukuje węzeł
o podanej wartości i zapamiętuje jego rodzica.
Jeżeli wartość nie występuje, drzewo pozostaje bez zmian.

Po odnalezieniu węzła rozpatrywane są trzy przypadki:
\begin{itemize}
    \item Brak dzieci: połączenie od rodzica zostaje usunięte.
    \item Jedno dziecko: rodzic zostaje połączony z tym dzieckiem.
    \item Dwoje dzieci: wartość węzła zostaje zastąpiona minimum
    prawego poddrzewa, po czym usuwany jest węzeł następnika.
\end{itemize}

Usuwanie korzenia wymaga odpowiedniej aktualizacji
pola \texttt{rootIndex}. Fizyczne usunięcie elementu
przez \texttt{vector::erase} przesuwa późniejsze elementy.
Dlatego wszystkie indeksy większe od usuniętej pozycji
są zmniejszane o jeden.

Udane usunięcie ma złożoność czasową $O(n)$ ze względu
na przesuwanie elementów i poprawianie indeksów.
Przy braku wartości wyszukiwanie zajmuje $O(h)$,
gdzie $h$ oznacza wysokość drzewa.
Dodatkowa pamięć wynosi $O(1)$.

\section{Usuwanie całego drzewa}

Metoda \texttt{removeTree()} wywołuje \texttt{tree.clear()}
oraz ustawia \texttt{rootIndex} na \texttt{-1}.
Po jej wykonaniu drzewo jest puste i gotowe
do ponownego dodawania elementów.

Metodę można wywoływać wielokrotnie.
Wektor zachowuje pojemność do ponownego wykorzystania.
Złożoność czasowa operacji wynosi $O(n)$,
a dodatkowa pamięć $O(1)$.

\section{Przykładowe sprawdzenie}

Program w funkcji \texttt{main} dodaje wartości
1, 2, 3, 8 oraz 5. Następnie usuwa wartość 3,
czyści całe drzewo i ponownie dodaje wartość 42.
Przykład sprawdza usunięcie węzła z jednym dzieckiem
oraz ponowne wykorzystanie wyczyszczonego obiektu.

Metoda \texttt{show()} wypisuje zawartość wektora.
Nie realizuje przejść preorder, inorder ani postorder.

\end{document}